\subsection{Composición constitutiva y evaluación de la sección gravitatoria}
\label{g26:subsec:evaluacion-viii-es}

La evaluación reutiliza las coordenadas regionales generadas y el registro
\(K\) seleccionado en el núcleo precedente. La coordenada \(\alpha\)
determina las secciones de acción de
\eqref{v:ant:eq:el-secciones-accion}; en la sección de retorno escribimos
\(\eta=\eta_{\rm ret}\) y \(S_*=10^{-34}\,\mathrm{J\,s}\), de modo
que \(\hbar_{\rm ret}=S_*\eta\). El ángulo y el contraángulo ya
construidos proporcionan, en la representación en grados,
\[
 A_{\deg}=1000\alpha,\qquad C^*_{\deg}=2(\eta+\alpha),\qquad
 x=\frac{\pi_{\rm HMT}A_{\deg}}{180},\quad
 y=\frac{\pi_{\rm HMT}C^*_{\deg}}{180}.
\]
Así, \(\eta=(C^*_{\deg}-2\alpha)/2\). La diferencia se toma entre
coordenadas de la misma representación angular.

En el dominio generado \(0<y<x\), los canales etiquetados son
\[
 q_\pm=e^{-(x\pm y)},\qquad
 r_\pm=\frac{1-q_\pm^{90}}{1-q_\pm^{120}},\qquad
 0<q_\pm<1.
\]
El cociente procede de los dos sectores de extensión: si
\(S_n(q)=\sum_{j=0}^{n-1}q^j\), la identidad
\((1-q)S_n(q)=1-q^n\) demuestra
\(r_\pm=S_{90}(q_\pm)/S_{120}(q_\pm)\).
Los cuadrados de los canales determinan las lecturas constitutivas
\[
 \widehat\varepsilon=r_+^2,\qquad
 \widehat\mu=r_-^2,\qquad
 \widehat c=\frac1{r_+r_-},\qquad
 \widehat Z=\frac{r_-}{r_+}.
\]
Para bases positivas de velocidad \(c_*\), acción \(u_S\) y carga
\(u_Q\), las secciones dimensionales se expresan como
\[
 \mu=\widehat\mu\frac{u_S}{c_*u_Q^2},\qquad
 \varepsilon=\widehat\varepsilon\frac{u_Q^2}{u_Sc_*},\qquad
 c=c_*\widehat c=\frac1{\sqrt{\mu\varepsilon}}.
\]
En efecto, el producto de las dos primeras ecuaciones cancela las bases
de acción y carga y deja
\(\mu\varepsilon=(r_+r_-)^2/c_*^2\); su raíz positiva prueba la
última igualdad. El factor \(\widehat c\) transforma la base \(c_*\)
una sola vez. La evaluación SI que sigue representa el mismo \(c\)
mediante \(299792458\,\mathrm{m\,s^{-1}}\) y restituye
\(c_*=c/\widehat c\). Esa operación verifica la equivalencia de las
representaciones dimensional y constitutiva.

Con \(r_N=5759/23040\) y la referencia declarada \(L_*=1\,\mathrm m\),
la sustitución de las ecuaciones anteriores en el coeficiente radial da
\begin{align}
 G&=\frac{c^3L_*^2}{S_*}\frac{r_N^2\alpha^{32}}{\eta}
   =\frac{2c^3L_*^2r_N^2\alpha^{32}}
           {S_*(C^*_{\deg}-2\alpha)},\label{g26:eq:g-angular-viii-es}\\
  &=\frac{L_*^2r_N^2\alpha^{32}}
           {S_*\eta(\mu\varepsilon)^{3/2}}.
 \label{g26:eq:g-constitutivo-viii-es}
\end{align}
Son expresiones del mismo coeficiente en la realización radial. La
referencia longitudinal \(L_*\), la longitud construida \(L_\alpha\),
la arista cronológica \(ct_0\) y el radio de horizonte conservan sus
tipos diferenciados.

La evaluación con las expansiones regionales archivadas proporciona
\begin{align*}
 L_\alpha&=1.6162549826251263301682625013\ldots\,10^{-35}\,\mathrm m,\\
 \hbar_{\rm ret}&=1.0545718169150390611336905847\ldots\,10^{-34}\,
                    \mathrm{J\,s},\\
 t_0&=3.1364999625365490734\ldots\,10^{-45}\,\mathrm s,\\
 G&=6.6742996594140227063677761891\ldots\,10^{-11}\,
                    \mathrm{m^3\,kg^{-1}\,s^{-2}}.
\end{align*}
Los cálculos a 100 y 140 cifras de trabajo conservan 71 cifras
significativas de la evaluación. Esta estabilidad numérica se distingue
de la incertidumbre experimental. Frente a la referencia CODATA 2022
\(G_{\rm ref}=6.67430(15)\,10^{-11}\,\mathrm{m^3\,kg^{-1}\,s^{-2}}\),
la diferencia es
\(G-G_{\rm ref}=-3.4058597729\ldots\,10^{-18}\,
\mathrm{m^3\,kg^{-1}\,s^{-2}}\), aproximadamente
\(-0.00227057\) veces su incertidumbre estándar. La referencia se
compara después del cálculo; las reglas R4, R5 y la identificación R8
conservan el alcance demostrado en la sección precedente.

Las fuentes de reproducción incluyen las expansiones regionales, el
registro de \(K\), las ecuaciones de acción y contraángulo, los
verificadores exactos y la evaluación decimal. Se ejecutan desde la
carpeta local \path{reproduccion_gravitatoria_20260926} mediante
\path{verificar_reproduccion.py}. Cada prueba declara su dominio; la
reutilización de coordenadas archivadas se distingue de una nueva
ejecución íntegra del generador.
